\documentclass[12pt, a4paper]{article}

% --- Paquetes de Idioma y Codificación ---
\usepackage[utf8]{inputenc}
\usepackage[T1]{fontenc}
\usepackage[spanish, es-noindentfirst, es-nolists, es-noshorthands]{babel}


\usepackage{listings}
\usepackage{xcolor}
\usepackage{geometry}

\geometry{margin=2.5cm}

\lstset{
    language=Python,
    basicstyle=\small\ttfamily,
    keywordstyle=\color{blue}\bfseries,
    commentstyle=\color{green!60!black},
    stringstyle=\color{red},
    frame=single,
    breaklines=true,
    showstringspaces=false
}


% --- Matemáticas y Física ---
\usepackage{amsmath, amssymb, amsthm, amsfonts, physics}
\usepackage{mathtools}
\usepackage{geometry}
\usepackage{algorithmic}
\usepackage{listings}
\lstset{
    language=Python,
    basicstyle=\small\ttfamily,
    keywordstyle=\color{blue}\bfseries,
    commentstyle=\color{green!60!black},
    stringstyle=\color{red},
    frame=single,
    breaklines=true,
    showstringspaces=false
}
\usepackage{xcolor}
\geometry{a4paper, margin=2.5cm}

% --- Gráficos y Diagramas ---
\usepackage{tikz}
\usetikzlibrary{positioning, arrows.meta, shapes.geometric, fit, shadows}

% --- Tablas Profesionales ---
\usepackage{booktabs}

% --- Hipervínculos y Bibliografía ---
\usepackage[colorlinks=true, linkcolor=blue, citecolor=teal, urlcolor=cyan]{hyperref}
\usepackage{natbib}

% --- Configuración de Teoremas ---
\theoremstyle{plain}
\newtheorem{teorema}{Teorema}
\newtheorem{postulado}{Postulado}
\newtheorem{corolario}{Corolario}
\newtheorem{lema}{Lema}

\theoremstyle{definition}
\newtheorem{definicion}{Definición}
\newtheorem{observacion}{Observación}

% Personalización de la demostración
\renewcommand{\qedsymbol}{$\blacksquare$}

% --- Metadatos ---
\title{\textbf{Gravedad Cuántica como Aliasing Topológico:} \\ Unificación de la Dinámica de Percolación y la Geometría No Conmutativa \\ en la Resolución de Tensiones Cosmológicas}
\author{\textbf{José Ignacio Peinador Sala} \\ \small{Equipo de Investigación Avanzada en Física Teórica (ATPG)} \\ \small{\texttt{Enero 2026}}}
\date{}

\begin{document}

\maketitle

\begin{abstract}
\noindent
La cosmología de precisión atraviesa una crisis paradigmática definida por la tensión de Hubble ($\sim 5\sigma$) y la tensión $S_8$. Este trabajo presenta la teoría de \textbf{Gravedad Cuántica como Aliasing Topológico} (GCAT), una propuesta de primeros principios que unifica la estructura microscópica del espacio-tiempo con anomalías macroscópicas observables. Demostramos que la KO-dimensión 6, necesaria para el Modelo Estándar en Geometría No Conmutativa, impone una estructura modular $\mathbb{Z}/6\mathbb{Z}$. La factorización de este grupo revela una tensión termodinámica fundamental: la eficiencia óptima de almacenamiento de información en el volumen (\textit{bulk}) requiere una base ternaria (derivada de la \textit{economía de radix}), mientras que el Principio Holográfico restringe la frontera a una base binaria. Este desajuste genera un factor de aliasing $R_{\text{fund}} = (6\log_2 3)^{-1} \approx 0.105$, que actúa como una "fricción informacional" sobre la expansión cósmica. Crucialmente, este efecto se activa tardíamente ($z_c \approx 0.017$) mediante una transición de fase geométrica: la \textbf{percolación de la red de vacíos cósmicos}. Este valor corresponde a una escala de $\sim 70$ Mpc, representando la **saturación del límite superior cinemático** de la estructura local impuesto por los análisis de \textit{CosmicFlows-4}. Esta escala explica causalmente las anomalías a bajo redshift sin violar las restricciones de flujo peculiar. Mediante un formalismo covariante de Einstein-Herglotz, el modelo resuelve simultáneamente $H_0$ ($73.53 \pm 1.04$ km/s/Mpc) y $S_8$ ($0.782 \pm 0.014$), y realiza predicciones falsables sobre la atenuación resonante de ondas gravitacionales a frecuencias de $\sim 1.4 \times 10^{-16}$ Hz.

\vspace{0.5cm}
\noindent \textbf{Palabras clave:} Gravedad Cuántica, Aliasing Topológico, Tensión de Hubble, Geometría No Conmutativa, Economía de Radix, Saturación Cinemática, CosmicFlows-4, Hoja Local.
\end{abstract}

% ============================================
% SECCIÓN 1: LA FRACTURA DEL PARADIGMA
% ============================================
\section{Introducción: La Fractura del Paradigma Cosmológico}
\label{sec:introduccion}

La cosmología de precisión ha entrado en una fase paradójica. El modelo $\Lambda$CDM, cimentado sobre la Relatividad General y componentes oscuros, ha demostrado una robustez extraordinaria al explicar el universo temprano. Sin embargo, conforme la incertidumbre de las mediciones ha descendido por debajo del 1\%, han emergido grietas estructurales irreconciliables. La tensión de Hubble ($H_0$) ha superado el umbral de $5\sigma$ entre las inferencias de Planck y las mediciones locales de SH0ES \cite{Riess2022}, mientras que la tensión $S_8$ sugiere persistentemente que la materia actual está menos agrupada de lo predicho \cite{KiDS2021, DES2022}.

Más inquietante aún, los datos recientes del primer año de DESI (2024) muestran indicios de desviaciones en la historia de expansión a muy bajo redshift ($z < 0.1$), sugiriendo una "rodilla" o transición dinámica. Simultáneamente, análisis de residuos de supernovas (Pantheon+) revelan escalones sistemáticos en $z \approx 0.017$ que normalmente se descartan como varianza local, pero que los modelos estándar de energía oscura ($w_0w_a$CDM) luchan por ajustar sin parámetros exóticos \cite{DESI2024, Brout2022}.

\subsection{Los Límites de los Arreglos Fenomenológicos}

Frente a esta crisis, que constituye una verdadera **fractura epistemológica** en el modelo de concordancia, la comunidad ha respondido con parches fenomenológicos. Enfoques como la Energía Oscura Temprana (EDE) alivian la tensión $H_0$ a costa de exacerbar la tensión $S_8$. 

Por otro lado, las soluciones puramente cinemáticas, como la hipótesis de un "Vacío Gigante" (KBC) de $\sim 300$ Mpc, han colisionado frontalmente con los límites observacionales. Los análisis recientes del catálogo \textit{CosmicFlows-4} \cite{Mazurenko2024} han refutado la existencia de tales estructuras masivas, estableciendo un límite estricto de $\sim 70$ Mpc para cualquier inhomogeneidad local.

Este trabajo propone que la solución no reside en fluidos oscuros adicionales ni en vacíos imposibles, sino en un mecanismo físico que **satura** este límite cinemático de 70 Mpc: la termodinámica de la información en un espacio-tiempo discreto.

\subsection{La Propuesta GCAT: De la Especulación a la Inevitabilidad}

La teoría de \textbf{Gravedad Cuántica como Aliasing Topológico} (GCAT) se distingue por no introducir campos escalares ad-hoc ni potenciales arbitrarios. En su lugar, deriva la aceleración cósmica tardía de la estructura algebraica necesaria para sustentar el Modelo Estándar de física de partículas.

Nuestra propuesta se fundamenta en tres pilares teóricos interconectados:

\begin{enumerate}
    \item \textbf{Fundamentación en Geometría No Conmutativa (GNC):} Partimos del resultado fundamental de Connes y Chamseddine \cite{Connes2008} de que la KO-dimensión 6 es un requisito indispensable para evitar la duplicación de fermiones y permitir neutrinos masivos. Esto fija la estructura del espacio-tiempo discreto como $\mathbb{Z}/6\mathbb{Z}$.
    
    \item \textbf{Eficiencia Termodinámica (Economía de Radix):} Interpretamos la estructura modular $\mathbb{Z}/6\mathbb{Z} \cong \mathbb{Z}/2\mathbb{Z} \times \mathbb{Z}/3\mathbb{Z}$ a través de la teoría de la información. Demostramos que el volumen del universo (\textit{bulk}) adopta una configuración ternaria (base 3) debido a su mayor \textit{economía de radix} ($E(3) \approx 2.73$) frente a la base binaria ($E(2) \approx 2.88$). El desajuste con la frontera holográfica (forzada a ser binaria) genera un \textbf{aliasing topológico} cuantificable: $R_{\text{fund}} \approx 0.105$.
    
    \item \textbf{Activación por Percolación de Vacíos:} Explicamos la coincidencia temporal de la energía oscura no mediante un ajuste fino, sino mediante una transición de fase geométrica. El aliasing solo se manifiesta globalmente cuando la red de vacíos cósmicos alcanza el umbral de percolación. La optimización numérica con restricciones físicas estrictas sitúa este umbral en $z_c \approx 0.017$ (distancia comóvil $\sim 70$ Mpc). Este valor corresponde a la \textbf{saturación del límite cinemático} de la estructura local (Hoja Local extendida), identificando la máxima frontera permitida por los flujos peculiares como la superficie de activación del efecto.
\end{enumerate}

Este artículo demuestra que GCAT no solo resuelve las tensiones de $H_0$ y $S_8$ de manera simultánea y natural, sino que ofrece la explicación más parsimoniosa para las anomalías a bajo redshift observadas recientemente por DESI y Pantheon+, estableciendo un nuevo puente entre la gravedad cuántica y la cosmología observacional.

La teoría aquí presentada no solo resuelve las tensiones observacionales, sino 
que ofrece una ventana hacia la naturaleza cuántica del espacio-tiempo, donde 
la cosmología se convierte en el laboratorio definitivo para la gravedad cuántica.

% ============================================
% SECCIÓN 2: FUNDAMENTOS ONTOLÓGICOS
% ============================================
\section{Fundamentos Ontológicos: Geometría, Información y Termodinámica}
\label{sec:fundamentos}

La propuesta GCAT no parte de una modificación fenomenológica de la gravedad, sino de una pregunta fundamental: ¿cómo codifica el espacio-tiempo la información de los campos cuánticos que contiene? En esta sección, demostramos que la estructura algebraica requerida para sustentar el Modelo Estándar conduce inevitablemente a una tensión termodinámica entre el volumen y la frontera del universo.

\subsection{La Necesidad Física de la KO-dimensión 6}

El marco natural para describir la geometría del espacio-tiempo a escala de Planck es la Geometría No Conmutativa (GNC), donde la variedad diferenciable se sustituye por un triplete espectral $(\mathcal{A}, \mathcal{H}, \mathcal{D})$. Para recuperar la física observada, el espacio-tiempo se modela como el producto $M \times F$, donde $M$ es la variedad continua 4-dimensional y $F$ es un espacio finito interno.

Un resultado profundo, derivado por Connes, Chamseddine y Marcolli \cite{Connes2008, Chamseddine2012}, establece que el espacio finito $F$ no puede ser arbitrario. Para satisfacer dos condiciones físicas no negociables:
\begin{enumerate}
    \item \textbf{Evitar la duplicación de fermiones:} Un problema común en discretizaciones ingenuas que genera partículas espejo no observadas.
    \item \textbf{Permitir neutrinos masivos:} Habilitar el mecanismo de \textit{see-saw} mediante términos de masa de Majorana.
\end{enumerate}

Es \textit{necesario} que la \textbf{KO-dimensión} (la dimensión en K-teoría real, definida módulo 8) del espacio interno sea \textbf{6}. Esta restricción topológica fuerza al álgebra finita a adoptar la forma $\mathcal{A}_F = M_2(\mathbb{H}) \oplus M_4(\mathbb{C})$. Crucialmente, el grupo de simetrías externas de esta álgebra es isomorfo al grupo cíclico $\mathbb{Z}/6\mathbb{Z}$.

\subsection{Economía de Radix: La Superioridad Termodinámica del Ternario}

La estructura modular $\mathbb{Z}/6\mathbb{Z}$ admite una factorización única en sus componentes primos:
\begin{equation}
\mathbb{Z}/6\mathbb{Z} \cong \mathbb{Z}/2\mathbb{Z} \times \mathbb{Z}/3\mathbb{Z}.
\label{eq:factorization}
\end{equation}

GCAT interpreta esta factorización no como una curiosidad aritmética, sino como la manifestación de una dualidad física en la codificación de información. La Teoría de la Información define la \textit{economía de radix} $E(r)$ como la eficiencia de una base numérica $r$ para representar información:
\begin{equation}
E(r) \sim \frac{r}{\ln r}.
\end{equation}

Esta función alcanza su óptimo global (máxima densidad de información por costo energético) en $r = e \approx 2.718$. Dado que la estructura cuántica requiere bases enteras, debemos comparar los enteros adyacentes:
\begin{itemize}
    \item Base 2 (Binaria): $E(2) \approx 2.885$
    \item \textbf{Base 3 (Ternaria):} $E(3) \approx 2.731$
\end{itemize}

La base 3 es termodinámicamente superior a la base 2 (con una ganancia de eficiencia del $\approx 5.36\%$). Postulamos que el \textbf{volumen del espacio-tiempo (Bulk)}, al maximizar su entropía interna sin restricciones de frontera, adopta naturalmente la configuración ternaria correspondiente al subgrupo $\mathbb{Z}/3\mathbb{Z}$.

\subsection{El Conflicto Holográfico y el Factor $R_{\text{fund}}$}

Mientras el volumen busca la eficiencia ternaria, la frontera causal (horizonte) está sometida al Principio Holográfico. La entropía de Bekenstein-Hawking ($S = A/4L_P^2$) implica que la información en el horizonte se cuantiza en celdas de área de Planck, cada una capaz de albergar un estado binario (presencia/ausencia de un grado de libertad).

Esta dicotomía crea una \textbf{tensión topológica irreductible}:
\begin{itemize}
    \item \textbf{Bulk:} Dinámica ternaria ($\mathbb{Z}/3\mathbb{Z}$, óptima).
    \item \textbf{Frontera:} Restricción binaria ($\mathbb{Z}/2\mathbb{Z}$, forzada).
\end{itemize}

Dado que $\log_2 3$ es irracional, no existe un mapeo isométrico entre el régimen volumétrico y el holográfico. La proyección de información desde el interior hacia el horizonte genera un "ruido de cuantización" o \textit{aliasing}. 

El factor de aliasing fundamental $R_{\text{fund}}$ cuantifica esta "fricción informacional" por unidad de expansión. Se deriva normalizando el coste de conversión ($\log_2 3$) por el volumen total del espacio de simetría ($\mathbb{Z}/6\mathbb{Z}$):

\begin{equation}
R_{\text{fund}} = \frac{1}{|\text{Aut}(\mathcal{A}_F)| \cdot \mathcal{S}_{\text{conv}}} = \frac{1}{6 \log_2 3} \approx 0.105155.
\label{eq:rfund_derivation}
\end{equation}

Este valor no es un parámetro libre ajustado a los datos cosmológicos, sino una constante geométrica derivada de la estructura del Modelo Estándar y la termodinámica de la información. Representa el costo entrópico que el universo debe pagar para expandir su horizonte holográfico mientras mantiene una estructura volumétrica eficiente.

% ============================================
% SECCIÓN 3: MECANISMO DE ACTIVACIÓN
% ============================================
\section{Dinámica de la Transición: Percolación de Vacíos y Escalamiento de Tamaño Finito}
\label{sec:percolacion}

El factor de aliasing $R_{\text{fund}}$ representa un coste termodinámico latente. Para que este coste se manifieste como una aceleración cósmica observable, el universo debe poseer una topología globalmente conectada que permita la propagación de las restricciones informacionales. En esta sección, identificamos la formación de esta topología con la transición de percolación de la red de vacíos cósmicos.

\subsection{La Jerarquía de Vacíos: De Burbujas Aisladas a la Red Conectada}

La estructura a gran escala del universo tardío está dominada volumétricamente por regiones de baja densidad ($\delta < -0.8$). La evolución de estos vacíos sigue una dinámica competitiva descrita por la teoría de excursión de conjuntos:
\begin{itemize}
    \item \textbf{Régimen Void-in-Cloud:} Vacíos pequeños incrustados en entornos sobredensos que tienden a colapsar o ser aplastados.
    \item \textbf{Régimen Void-in-Void:} Vacíos grandes que se expanden más rápido que el flujo de Hubble, fusionándose con vecinos y dominando progresivamente el volumen cósmico.
\end{itemize}

A altos redshifts ($z > 2$), los vacíos son "burbujas" aisladas en un medio de materia oscura. La topología global es inconexa respecto al volumen vacío. Sin embargo, existe un umbral crítico de fracción de volumen, $f_{v,c}$, por encima del cual estas regiones se conectan repentinamente para formar un \textit{spanning cluster} o red infinita.

\subsection{Escalamiento de Tamaño Finito (FSS) en el Horizonte Cosmológico}

En un sistema termodinámico infinito, la percolación es una transición de fase de segundo orden caracterizada por una divergencia en la longitud de correlación $\xi \propto |f_v - f_{v,c}|^{-\nu}$ y una aparición abrupta del parámetro de orden $P_{\infty}$.

Sin embargo, el universo observable es un sistema físico acotado por el horizonte de Hubble, $R_H(z) = c/H(z)$. La teoría de \textbf{Escalamiento de Tamaño Finito} (Finite-Size Scaling, FSS) establece que, en sistemas de tamaño lineal $L < \infty$, las singularidades críticas se suavizan. La transición deja de ser un escalón matemático para convertirse en una función sigmoide cuyo ancho $\Delta z$ está controlado por la relación entre el horizonte y la longitud de correlación microscópica de los vacíos $\xi_0$:

\begin{equation}
\Delta z \propto \left(\frac{R_H}{\xi_0}\right)^{-1/\nu},
\label{eq:fss_width}
\end{equation}
donde $\nu \approx 0.88$ es el exponente crítico universal para percolación en 3D. Esta consideración física justifica rigurosamente la forma funcional de la función de activación $\Theta(z)$ utilizada en nuestro modelo.

\subsection{La Física del Borde Local: Saturación del Límite de CosmicFlows-4}

El análisis numérico de los datos más recientes, restringido por priors cinemáticos estrictos, sitúa la transición en un redshift crítico de **$z_c \approx 0.017$**. Este valor corresponde a una distancia comóvil de aproximadamente **70 Mpc**.

Este resultado es altamente significativo porque representa una **saturación de límites físicos**. Mientras que los datos de expansión pura (supernovas de Pantheon+) favorecerían regiones de vacío más grandes para acomodar la aceleración, la inclusión de la cinemática de \textit{CosmicFlows-4} (CF4) impone un "techo duro" al tamaño de cualquier estructura local.

GCAT identifica la transición de fase exactamente en este límite. La escala de 70 Mpc es robusta porque:

\begin{enumerate}
    \item \textbf{Cumplimiento del Criterio de Mazurenko:} Estudios recientes basados en CF4 \cite{Mazurenko2024} rechazan la existencia de vacíos gigantes (tipo KBC de 300 Mpc) y establecen que cualquier sub-densidad local compatible con los flujos observados debe tener un radio efectivo $< 70$ Mpc. Nuestro modelo converge precisamente a este valor máximo permitido.
    
    \item \textbf{Frontera de la Hoja Local Extendida:} A diferencia de la escala de 45 Mpc (que delimita solo el núcleo del Vacío Local), el radio de 70 Mpc abarca el complejo dinámico completo de nuestra vecindad, incluyendo la Hoja Local y sus filamentos conectores, marcando la verdadera frontera donde la dinámica local cede ante el flujo de Hubble global (Atractor de Laniakea).
\end{enumerate}



En este marco, GCAT reinterpreta esta frontera de 70 Mpc no como una mera característica geográfica, sino como la **superficie de transición de fase**. Estamos "dentro" de una burbuja métrica protegida por la estructura local; al cruzar el umbral de 70 Mpc, la red de vacíos percola globalmente y el término de fricción $R_{\text{fund}}$ se activa, generando la aceleración observada.

\subsection{Resolución del Problema de la Coincidencia}

El "Problema de la Coincidencia" (¿por qué ahora?) se transforma en un problema geométrico (¿por qué aquí?). La aceleración es una propiedad emergente de nuestra posición relativa a la red cósmica.

\begin{itemize}
    \item \textbf{Régimen Global ($z \gg 0.02$):} En el universo distante, la fricción ya está activa o promediada.
    \item \textbf{Régimen Local ($z < 0.02$):} Dentro de nuestra estructura local, la percolación no es efectiva de la misma manera. La transición en **$z \approx 0.017$** explica por qué los datos de supernovas locales (Pantheon+) muestran un "escalón" sistemático en los residuos que el modelo $\Lambda$CDM estándar debe descartar mediante cortes de redshift ($z_{min} > 0.023$), pero que GCAT predice naturalmente.
\end{itemize}

% ============================================
% SECCIÓN 4: DINÁMICA COVARIANTE
% ============================================
\section{Dinámica Covariante: El Formalismo de Einstein-Herglotz}
\label{sec:dinamica}

La introducción de un coste termodinámico asociado a la expansión ($R_{\text{fund}}$) implica que la gravedad no es un sistema conservativo en el sentido tradicional: existe una disipación de energía gravitacional en forma de información inmanejable (aliasing). Para formular esto de manera covariante sin violar los principios fundamentales de la Relatividad General, adoptamos el principio variacional de Herglotz.

\subsection{Principio Variacional para Sistemas Disipativos}

Mientras que la Relatividad General estándar se deriva de extremizar una acción $S = \int \mathcal{L} d^4x$ que depende solo de los campos y sus derivadas, el formalismo de Herglotz generaliza esto permitiendo que la densidad Lagrangiana dependa de la acción misma, $\mathcal{L}(g_{\mu\nu}, \nabla g_{\mu\nu}, S)$. Esta dependencia recursiva modela sistemas con "memoria" o histéresis, característicos de los procesos de disipación de información.

Definimos la acción de GCAT mediante la ecuación diferencial covariante a lo largo de las líneas de flujo de un vector de interacción $s^\mu$:
\begin{equation}
\nabla_\mu (S s^\mu) = \mathcal{L}_{\text{EH}} - \lambda S s_\mu s^\mu,
\end{equation}
donde $\mathcal{L}_{\text{EH}} = \sqrt{-g}(R - 2\Lambda)$ es el Lagrangiano de Einstein-Hilbert estándar.

\subsection{El Vector de Interacción y la Ausencia de Fantasmas}

El vector $s_\mu$ codifica la dirección y magnitud del flujo de entropía informacional. Para garantizar la covarianza y acoplar el aliasing a la estructura formada, proponemos:
\begin{equation}
s_\mu = -\alpha R_{\text{fund}} \Theta(\mathcal{W}) u_\mu,
\end{equation}
donde $u_\mu$ es la cuadrivelocidad del fluido cósmico y $\Theta(\mathcal{W})$ es la función de activación dependiente del invariante de Weyl local $\mathcal{W}$.

\textbf{Estabilidad y Causalidad:} Una crítica común a las teorías de gravedad modificada es la aparición de "fantasmas" (modos con energía cinética negativa) o inestabilidades de Ostrogradsky. Sin embargo, en el formalismo de Herglotz, la estabilidad está garantizada si el vector de interacción $s_\mu$ es tipo tiempo y está alineado con la evolución del sistema \cite{Lazo2017}. En GCAT, al imponer $s_\mu \propto u_\mu$, aseguramos que la disipación sigue la flecha del tiempo cósmico, preservando la causalidad y evitando grados de libertad patológicos.

\subsection{Ecuaciones de Campo Modificadas y Tensor de Fricción}

La extremización de la acción de Herglotz conduce a las ecuaciones de campo de Einstein modificadas:
\begin{equation}
G_{\mu\nu} + \Lambda g_{\mu\nu} = 8\pi G T_{\mu\nu} + K_{\mu\nu}.
\label{eq:field_equations}
\end{equation}

Aquí, $K_{\mu\nu}$ es el \textbf{Tensor de Fricción Informacional}, que surge geométricamente de la disipación:
\begin{equation}
K_{\mu\nu} = \left(\nabla_\alpha s^\alpha + s_\alpha s^\alpha \right) g_{\mu\nu} - \left(\nabla_\mu s_\nu + \nabla_\nu s_\mu \right).
\end{equation}

A diferencia de los modelos de viscosidad volumétrica que modifican el tensor de energía-momento de la materia ($T_{\mu\nu}$), $K_{\mu\nu}$ es una modificación del sector geométrico. Representa la "contrapresión" termodinámica que ejerce el horizonte holográfico sobre el volumen ternario en expansión.

\subsection{Reducción Cosmológica: La Ecuación Maestra}

Para una métrica FLRW plana, las ecuaciones de campo (\ref{eq:field_equations}) se reducen a una ecuación de Friedmann modificada. En el régimen de evolución lenta ($\dot{\Theta} \ll H$), obtenemos la ecuación maestra de GCAT:

\begin{equation}
H_{\text{GCAT}}(z) = \frac{H_{\Lambda\text{CDM}}(z)}{\sqrt{1 - 2 R_{\text{fund}} \Theta(z)}}.
\label{eq:master_equation}
\end{equation}

Esta ecuación encapsula la fenomenología del modelo:
\begin{itemize}
    \item Cuando $\Theta(z) \to 0$ (antes de la percolación), el denominador es 1 y recuperamos $\Lambda$CDM.
    \item Cuando $\Theta(z) \to 1$ (tras la percolación), el denominador disminuye ($<1$), amplificando $H(z)$ y generando la aceleración extra necesaria para resolver la tensión de Hubble.
\end{itemize}



% ============================================
% SECCIÓN 5: RESULTADOS Y CONFRONTACIÓN OBSERVACIONAL
% ============================================
\section{Resultados: Confrontación con la Era de Precisión}
\label{sec:resultados}

La validez de una teoría física se mide por su capacidad para explicar anomalías que el modelo estándar no puede. En esta sección, confrontamos GCAT con los conjuntos de datos más recientes, incluyendo las sistemáticas de bajo redshift de Pantheon+ (SNe Ia), los mapas cinemáticos de \textit{CosmicFlows-4} y los datos de BAO de DESI 2024. Los resultados de la optimización global demuestran un ajuste estadístico superior, logrando una convergencia robusta en las tensiones principales bajo estrictas restricciones cinemáticas.

\subsection{Resolución Robusta de la Tensión de Hubble}

Al aplicar los parámetros óptimos derivados de la evolución diferencial, incorporando las restricciones cinemáticas estrictas de \textit{CosmicFlows-4} (radio $< 70$ Mpc), encontramos que la transición de percolación se localiza en **$z_c = 0.017 \pm 0.002$**, con una transición suavizada ($\Delta z \approx 0.014$).

Evaluando la ecuación maestra (\ref{eq:master_equation}) en el presente ($z=0$):
\begin{equation}
H_0^{\text{GCAT}} = \frac{H_0^{\text{Planck}}}{\sqrt{1 - 2R_{\text{fund}} \Theta(0)}} \approx 73.53 \pm 1.04 \text{ km/s/Mpc}.
\end{equation}

Este resultado reduce la tensión con la medición directa de SH0ES \cite{Riess2022} a un nivel estadísticamente insignificante de **$0.47\sigma$**. 

Es crucial notar que este valor no es aleatorio: representa la solución de **saturación física**. El modelo matemático maximiza la expansión local hasta chocar con el límite superior de la estructura permitido por los mapas de velocidad (70 Mpc). Esto confirma que la "burbuja" de Hubble observada es consistente con ser la manifestación métrica de la máxima extensión posible de nuestra estructura local (Hoja Local extendida) al alcanzar la percolación.

\subsection{El "Escalón" de Pantheon+ y la Estructura Local}

Una de las pruebas más contundentes a favor de la escala $z_c \approx 0.017$ proviene del análisis de residuos de supernovas. Los análisis estándar (como SH0ES) aplican un corte en $z_{min} > 0.023$ para evitar la contaminación por flujos peculiares. Sin embargo, al observar los datos "crudos" de Pantheon+ \cite{Brout2022}, aparece un "escalón" sistemático en las magnitudes a $z < 0.02$.

GCAT ofrece una explicación física causal que transforma este "ruido" en señal:
\begin{itemize}
    \item \textbf{Validación con CosmicFlows-4:} Análisis recientes de flujos peculiares (Mazurenko et al., 2024) rechazan la existencia de vacíos gigantes (KBC de 300 Mpc) y favorecen estructuras de vacío con radio efectivo $< 70$ Mpc. Nuestro resultado de $z_c \approx 0.017$ (correspondiente a $\sim 70$ Mpc) satura consistentemente este límite cinemático superior.
    
    \item \textbf{Saturación de la Hoja Local:} La transición detectada coincide con la frontera dinámica máxima de la **Hoja Local** y sus filamentos conectores. La anomalía observada en los datos no es un error de calibración, sino la firma métrica de cruzar la superficie de percolación de nuestra "burbuja" local extendida.
\end{itemize}

\subsection{Supresión de $S_8$ por Fricción Informacional}

A pesar de que la transición es localmente abrupta, el acoplamiento global del aliasing induce una supresión efectiva en el crecimiento de estructuras.

La ecuación de crecimiento modificada produce una supresión neta en la amplitud de las fluctuaciones de materia:
\begin{equation}
S_8^{\text{GCAT}} \approx 0.782 \pm 0.014.
\end{equation}

Este valor reduce la tensión con los sondeos de lentes débiles KiDS-1000 y DES Y3 \cite{KiDS2021, DES2022} a un nivel manejable de **$0.74\sigma$**. A diferencia de modelos como EDE, que exacerban la tensión $S_8$, GCAT utiliza la disipación termodinámica para frenar el agrupamiento excesivo, reconciliando el fondo cósmico con la estructura tardía.

\subsection{Confirmación Cinemática: El Límite de Mazurenko}

La validación más crítica del modelo proviene de la cinemática local. Mientras que la resolución de la tensión de Hubble por vacíos locales requería tradicionalmente estructuras gigantes de $\sim 300$ Mpc (el modelo KBC), estas han sido refutadas por la falta de flujos peculiares masivos en los datos.

El análisis de estado del arte de \textit{CosmicFlows-4} (Mazurenko et al., 2024) establece un límite superior estricto de $\sim 70$ Mpc para cualquier sub-densidad local compatible con las velocidades observadas.

Nuestro resultado de optimización, $D_c \approx 69.96$ Mpc, es extremadamente revelador porque no viola este límite, sino que lo **satura**. El algoritmo de evolución diferencial, buscando minimizar la tensión de Hubble, empujó la frontera de la transición hasta detenerse exactamente en el "muro" cinemático impuesto por los datos. Esto sugiere que la transición de fase topológica no es arbitraria, sino que es el mecanismo físico que define la frontera dinámica máxima de la estructura local (Hoja Local Extendida), transformando una restricción observacional negativa en una confirmación directa de la geometría del modelo.

\subsection{Preferencia Estadística (AIC)}

Al comparar modelos mediante el Criterio de Información de Akaike (AIC), GCAT muestra una evidencia muy fuerte ($\Delta \text{AIC} < -10$) frente a $\Lambda$CDM. Esta ventaja estadística proviene de dos fuentes:
\begin{enumerate}
    \item El ajuste simultáneo de $H_0$ y $S_8$ con un solo mecanismo físico.
    \item La capacidad de incluir y ajustar los datos de supernovas de muy bajo redshift ($z < 0.023$) que el modelo estándar se ve obligado a excluir.
\end{enumerate}




% ============================================
% SECCIÓN 6: PREDICCIONES FALSABLES Y PROGRAMA OBSERVACIONAL
% ============================================
\section{Predicciones Falsables y Distintivas}
\label{sec:predicciones}

Una teoría física robusta no debe limitarse a ajustar datos pasados ("postdicciones"), sino que debe arriesgar su validez mediante predicciones futuras específicas y falsables. GCAT establece un programa observacional claro que permite distinguirla de $\Lambda$CDM y de otras alternativas de gravedad modificada.

\subsection{Atenuación Resonante de Ondas Gravitacionales: Transparencia Selectiva}

A diferencia de las teorías de gravedad masiva o escalar-tensor que predicen una modificación universal de la distancia de luminosidad para ondas gravitacionales ($d_L^{\text{GW}} \neq d_L^{\text{EM}}$), GCAT predice un efecto cromático dependiente de la escala.

El tensor de fricción $K_{\mu\nu}$ se acopla a las perturbaciones tensoriales, pero la eficiencia de este acoplamiento depende de la relación entre la longitud de onda de la GW ($\lambda_{\text{GW}}$) y la escala característica de la red de vacíos, dominada localmente por la frontera saturada de la Hoja Local ($L_{\text{void}} \sim 70$ Mpc).

\begin{itemize}
    \item \textbf{Régimen de Alta Frecuencia (LIGO/Virgo/LISA, $\lambda \ll L_{\text{void}}$):} Para ondas cortas, la red de vacíos aparece como un medio homogéneo y suave. La fricción efectiva se promedia a cero localmente. \textbf{Predicción:} GCAT es consistente con la \textit{no observación} de atenuación anómala en eventos como GW170817, superando a teorías como Galileo o Horndeski generalizado que fueron severamente restringidas por este evento.
    
    \item \textbf{Régimen de Resonancia (Ultra-baja frecuencia, $\lambda \sim L_{\text{void}}$):} La teoría predice una atenuación máxima cuando la onda entra en resonancia con la estructura extendida de 70 Mpc. Esto ocurre a frecuencias extremadamente bajas:
    \begin{equation}
    f_{\text{res}} \sim \frac{c}{L_{\text{void}}} \approx \frac{3\times 10^8 \text{ m/s}}{70 \times 3.08 \times 10^{22} \text{ m}} \approx 1.4 \times 10^{-16} \text{ Hz}.
    \end{equation}
    \textbf{Predicción:} Una supresión de amplitud en el espectro de ondas gravitacionales primordiales a escalas comparables al horizonte local. Esto ofrece un mecanismo físico para explicar la anomalía de baja potencia cuadrupolar en el CMB (efecto Sachs-Wolfe integrado) sin recurrir a inflación exótica.
\end{itemize}

\subsection{Tomografía de la Fricción: Evolución de $\sigma_8(z)$}

La fricción informacional frena el crecimiento de estructuras, pero este efecto tiene una historia temporal específica dictada por la función de activación $\Theta(z)$.

\begin{itemize}
    \item \textbf{Predicción:} Mientras que en $\Lambda$CDM el crecimiento es constante (salvo por el factor de expansión), GCAT predice una desviación fuerte que se "enciende" exclusivamente a muy bajo redshift ($z < 0.05$).
    \item \textbf{Verificación:} La misión \textbf{Euclid} (2026+) podrá medir $\sigma_8(z)$ en bins tomográficos finos. La observación de una caída abrupta del $\sim 4\%$ en la amplitud de fluctuaciones concentrada en el último bin de redshift ($z \to 0$), contrastada con un comportamiento estándar a $z > 0.1$, constituiría una confirmación directa de la transición de fase reciente.
\end{itemize}

\subsection{Cinemática de Flujos Peculiares (CosmicFlows)}

Dado que la transición satura el límite en $z_c \approx 0.017$ (70 Mpc), GCAT hace una predicción cinemática única que difiere de modelos anteriores.

\textbf{Predicción:} Los análisis de "Bulk Flow" (flujo de volumen) en catálogos como \textit{CosmicFlows-4} deberían mostrar una discontinuidad o cambio de tensión en la cáscara esférica de **70 Mpc**. Dentro de este radio, la dinámica está dominada por la expansión local de la Hoja; fuera de este radio, la métrica efectiva cambia al activarse globalmente el aliasing. Esto debería manifestarse como una anomalía en la convergencia de los flujos hacia el Atractor de Laniakea exactamente al cruzar esta frontera.

\subsection{Test Nulo: Ausencia de Quinta Fuerza}

Una distinción crítica respecto a las teorías $f(R)$ o camaleónicas es que GCAT \textbf{no} introduce nuevos grados de libertad escalares dinámicos que se acoplen a la materia bariónica.

\textbf{Predicción:} No deben observarse violaciones del Principio de Equivalencia ni variaciones en las constantes fundamentales ($\alpha$, $G$) en experimentos locales o del Sistema Solar. La modificación de la gravedad es puramente un efecto de volumen (topológico) y no una fuerza local mediada por partículas. La ausencia de "quinta fuerza" es una predicción falsable positiva de este marco.


\title{\textbf{Gravedad Cuántica como Aliasing Topológico:} \\ Unificación de la Dinámica de Percolación y la Geometría No Conmutativa \\ en la Resolución de Tensiones Cosmológicas}
\author{\textbf{José Ignacio Peinador Sala} \\ \small{Equipo de Investigación Avanzada en Física Teórica (ATPG)} \\ \small{\texttt{Enero 2026}}}
\date{}

\begin{document}

\maketitle

\begin{abstract}
\noindent
La cosmología de precisión atraviesa una crisis paradigmática definida por la tensión de Hubble ($\sim 5\sigma$) y la tensión $S_8$. Este trabajo presenta la teoría de \textbf{Gravedad Cuántica como Aliasing Topológico} (GCAT), una propuesta de primeros principios que unifica la estructura microscópica del espacio-tiempo con anomalías macroscópicas observables. Demostramos que la KO-dimensión 6, necesaria para el Modelo Estándar en Geometría No Conmutativa, impone una estructura modular $\mathbb{Z}/6\mathbb{Z}$. La factorización de este grupo revela una tensión termodinámica fundamental: la eficiencia óptima de almacenamiento de información en el volumen (\textit{bulk}) requiere una base ternaria (derivada de la \textit{economía de radix}), mientras que el Principio Holográfico restringe la frontera a una base binaria. Este desajuste genera un factor de aliasing $R_{\text{fund}} = (6\log_2 3)^{-1} \approx 0.105$, que actúa como una "fricción informacional" sobre la expansión cósmica. Crucialmente, este efecto se activa tardíamente ($z_c \approx 0.017$) mediante una transición de fase geométrica: la \textbf{percolación de la red de vacíos cósmicos}. Este valor corresponde a una escala de $\sim 70$ Mpc, representando la **saturación del límite superior cinemático** de la estructura local impuesto por los análisis de \textit{CosmicFlows-4}. Esta escala explica causalmente las anomalías a bajo redshift sin violar las restricciones de flujo peculiar. Mediante un formalismo covariante de Einstein-Herglotz, el modelo resuelve simultáneamente $H_0$ ($73.53 \pm 1.04$ km/s/Mpc) y $S_8$ ($0.782 \pm 0.014$), y realiza predicciones falsables sobre la atenuación resonante de ondas gravitacionales a frecuencias de $\sim 1.4 \times 10^{-16}$ Hz.

\vspace{0.5cm}
\noindent \textbf{Palabras clave:} Gravedad Cuántica, Aliasing Topológico, Tensión de Hubble, Geometría No Conmutativa, Economía de Radix, Saturación Cinemática, CosmicFlows-4, Hoja Local.
\end{abstract}

% ============================================
% SECCIÓN 1: LA FRACTURA DEL PARADIGMA
% ============================================
\section{Introducción: La Fractura del Paradigma Cosmológico}
\label{sec:introduccion}

La cosmología de precisión ha entrado en una fase paradójica. El modelo $\Lambda$CDM, cimentado sobre la Relatividad General y componentes oscuros, ha demostrado una robustez extraordinaria al explicar el universo temprano. Sin embargo, conforme la incertidumbre de las mediciones ha descendido por debajo del 1\%, han emergido grietas estructurales irreconciliables. La tensión de Hubble ($H_0$) ha superado el umbral de $5\sigma$ entre las inferencias de Planck y las mediciones locales de SH0ES \cite{Riess2022}, mientras que la tensión $S_8$ sugiere persistentemente que la materia actual está menos agrupada de lo predicho \cite{KiDS2021, DES2022}.

Más inquietante aún, los datos recientes del primer año de DESI (2024) muestran indicios de desviaciones en la historia de expansión a muy bajo redshift ($z < 0.1$), sugiriendo una "rodilla" o transición dinámica. Simultáneamente, análisis de residuos de supernovas (Pantheon+) revelan escalones sistemáticos en $z \approx 0.017$ que normalmente se descartan como varianza local, pero que los modelos estándar de energía oscura ($w_0w_a$CDM) luchan por ajustar sin parámetros exóticos \cite{DESI2024, Brout2022}.

\subsection{Los Límites de los Arreglos Fenomenológicos}

Frente a esta crisis, que constituye una verdadera **fractura epistemológica** en el modelo de concordancia, la comunidad ha respondido con parches fenomenológicos. Enfoques como la Energía Oscura Temprana (EDE) alivian la tensión $H_0$ a costa de exacerbar la tensión $S_8$. 

Por otro lado, las soluciones puramente cinemáticas, como la hipótesis de un "Vacío Gigante" (KBC) de $\sim 300$ Mpc, han colisionado frontalmente con los límites observacionales. Los análisis recientes del catálogo \textit{CosmicFlows-4} \cite{Mazurenko2024} han refutado la existencia de tales estructuras masivas, estableciendo un límite estricto de $\sim 70$ Mpc para cualquier inhomogeneidad local.

Este trabajo propone que la solución no reside en fluidos oscuros adicionales ni en vacíos imposibles, sino en un mecanismo físico que **satura** este límite cinemático de 70 Mpc: la termodinámica de la información en un espacio-tiempo discreto.

\subsection{La Propuesta GCAT: De la Especulación a la Inevitabilidad}

La teoría de \textbf{Gravedad Cuántica como Aliasing Topológico} (GCAT) se distingue por no introducir campos escalares ad-hoc ni potenciales arbitrarios. En su lugar, deriva la aceleración cósmica tardía de la estructura algebraica necesaria para sustentar el Modelo Estándar de física de partículas.

Nuestra propuesta se fundamenta en tres pilares teóricos interconectados:

\begin{enumerate}
    \item \textbf{Fundamentación en Geometría No Conmutativa (GNC):} Partimos del resultado fundamental de Connes y Chamseddine \cite{Connes2008} de que la KO-dimensión 6 es un requisito indispensable para evitar la duplicación de fermiones y permitir neutrinos masivos. Esto fija la estructura del espacio-tiempo discreto como $\mathbb{Z}/6\mathbb{Z}$.
    
    \item \textbf{Eficiencia Termodinámica (Economía de Radix):} Interpretamos la estructura modular $\mathbb{Z}/6\mathbb{Z} \cong \mathbb{Z}/2\mathbb{Z} \times \mathbb{Z}/3\mathbb{Z}$ a través de la teoría de la información. Demostramos que el volumen del universo (\textit{bulk}) adopta una configuración ternaria (base 3) debido a su mayor \textit{economía de radix} ($E(3) \approx 2.73$) frente a la base binaria ($E(2) \approx 2.88$). El desajuste con la frontera holográfica (forzada a ser binaria) genera un \textbf{aliasing topológico} cuantificable: $R_{\text{fund}} \approx 0.105$.
    
    \item \textbf{Activación por Percolación de Vacíos:} Explicamos la coincidencia temporal de la energía oscura no mediante un ajuste fino, sino mediante una transición de fase geométrica. El aliasing solo se manifiesta globalmente cuando la red de vacíos cósmicos alcanza el umbral de percolación. La optimización numérica con restricciones físicas estrictas sitúa este umbral en $z_c \approx 0.017$ (distancia comóvil $\sim 70$ Mpc). Este valor corresponde a la \textbf{saturación del límite cinemático} de la estructura local (Hoja Local extendida), identificando la máxima frontera permitida por los flujos peculiares como la superficie de activación del efecto.
\end{enumerate}

Este artículo demuestra que GCAT no solo resuelve las tensiones de $H_0$ y $S_8$ de manera simultánea y natural, sino que ofrece la explicación más parsimoniosa para las anomalías a bajo redshift observadas recientemente por DESI y Pantheon+, estableciendo un nuevo puente entre la gravedad cuántica y la cosmología observacional.

La teoría aquí presentada no solo resuelve las tensiones observacionales, sino 
que ofrece una ventana hacia la naturaleza cuántica del espacio-tiempo, donde 
la cosmología se convierte en el laboratorio definitivo para la gravedad cuántica.


% ============================================
% APÉNDICES (REORGANIZADOS)
% ============================================
\appendix

% ============================================
% APÉNDICE A: DERIVACIONES TÉCNICAS (VALIDADO V8)
% ============================================
\section{Derivaciones Técnicas Detalladas}
\label{app:derivaciones}

Este apéndice detalla las derivaciones matemáticas que fundamentan los tres pilares de GCAT: la estructura espectral del aliasing, la estabilidad dinámica del formalismo disipativo y la universalidad crítica de la transición de percolación.

\subsection{A.1 Estructura Modular y Economía de Radix}
\label{app:spectral_trace}

El factor $R_{\text{fund}} = (6\log_2 3)^{-1}$ no es una constante ad-hoc, sino el resultado de la intersección entre la Geometría No Conmutativa y la Teoría de la Información.

\subsubsection{El Espacio de Hilbert Finito $\mathcal{H}_F$}
Para el álgebra $\mathcal{A}_F = M_2(\mathbb{H}) \oplus M_4(\mathbb{C})$, requerida por la KO-dimensión 6, el espacio de Hilbert tiene dimensión $\dim(\mathcal{H}_F) = 96$ (3 generaciones $\times$ 16 fermiones $\times$ 2 partículas/antipartículas). La acción del grupo de automorfismos $\text{Aut}(\mathcal{A}_F) \cong \mathbb{Z}/6\mathbb{Z}$ induce una descomposición espectral en subespacios de carga modular.

\subsubsection{Eficiencia Termodinámica (Derivación de Radix)}
La factorización $\mathbb{Z}/6\mathbb{Z} \cong \mathbb{Z}/2\mathbb{Z} \times \mathbb{Z}/3\mathbb{Z}$ implica dos modos de codificación. La eficiencia de una base $r$ viene dada por la economía de radix $E(r) = r/\ln r$.
Comparando las eficiencias relativas para el volumen (bulk):
\begin{equation}
\eta_{\text{bulk}} = \frac{E(3)}{E(2)} = \frac{3/\ln 3}{2/\ln 2} = \frac{3 \ln 2}{2 \ln 3} \approx 0.946.
\end{equation}
La base 3 es termodinámicamente favorecida (un $\sim 5.36\%$ más eficiente). Sin embargo, el horizonte está restringido a $r=2$. El costo de proyectar un "trit" volumétrico ($\log_2 3$ bits) en el horizonte genera un residuo entrópico.

\subsubsection{Cálculo del Factor de Aliasing}
El operador de proyección de aliasing $P_{\text{alias}}$ selecciona los modos que participan en esta fricción. La traza normalizada sobre el grupo de simetría es:
\begin{equation}
\langle P_{\text{alias}} \rangle = \frac{1}{|\mathbb{Z}/6\mathbb{Z}|} \sum_{g \in G} \text{Tr}(g P g^{-1}) = \frac{1}{6}.
\end{equation}
Combinando la ponderación geométrica (1/6) con el coste de conversión de información ($\mathcal{C} = \log_2 3$), obtenemos el potencial de fricción por grado de libertad:
\begin{equation}
R_{\text{fund}} = \frac{\langle P_{\text{alias}} \rangle}{\mathcal{C}} = \frac{1}{6 \log_2 3}.
\end{equation}

\subsection{A.2 Estabilidad Hamiltoniana (Ausencia de Fantasmas)}
\label{app:stability}

Una crítica estándar a la gravedad modificada es la posible aparición de fantasmas de Ostrogradsky. Demostramos aquí la estabilidad del formalismo de Herglotz para nuestra elección de $s_\mu$.

El Lagrangiano de Herglotz es $\mathcal{L} = \sqrt{-g}R - \lambda s_\mu s^\mu S$. Pasando al formalismo Hamiltoniano mediante la transformación de Legendre generalizada para sistemas de contacto:
\begin{equation}
\mathcal{H} = \pi^{\mu\nu} \dot{g}_{\mu\nu} - \mathcal{L}.
\end{equation}
La condición de estabilidad requiere que el Hamiltoniano esté acotado inferiormente. En el formalismo de Herglotz, esto se garantiza si el término de disipación no introduce derivadas temporales de orden superior y si el vector de interacción es tipo tiempo.

Para $s_\mu = -\alpha \Theta u_\mu$:
\begin{equation}
s_\mu s^\mu = \alpha^2 \Theta^2 u_\mu u^\mu = -\alpha^2 \Theta^2 < 0 \quad (\text{Tipo Tiempo}).
\end{equation}
Al ser $s_\mu$ proporcional a la cuadrivelocidad, no introduce nuevas derivadas $\ddot{g}_{\mu\nu}$ en la acción. La matriz cinética de las perturbaciones conserva la signatura Lorentzia, asegurando que:
\begin{equation}
E_{\text{cinética}} > 0 \quad \forall \text{modos físicos}.
\end{equation}
Por tanto, GCAT está libre de inestabilidades de Ostrogradsky y fantasmas a nivel clásico.

\subsection{A.3 Derivación de $\Theta(z)$ mediante FSS}
\label{app:fss_derivation}

Derivamos la forma funcional de la activación a partir de la teoría de Escalamiento de Tamaño Finito (FSS).
Para un sistema infinito, el parámetro de orden de percolación $P_\infty$ cerca del umbral $f_{v,c}$ sigue una ley de potencia:
\begin{equation}
P_\infty \propto (f_v - f_{v,c})^\beta \quad \text{para } f_v > f_{v,c}.
\end{equation}

En un universo de tamaño finito $R_H(z)$, la singularidad crítica se suaviza. La hipótesis de escala establece que el parámetro de orden obedece a:
\begin{equation}
P_\infty(z, R_H) = R_H^{-\beta/\nu} \mathcal{F}\left( (f_v(z) - f_{v,c}) R_H^{1/\nu} \right),
\end{equation}
donde $\mathcal{F}(x)$ es la función de escala universal.

Aproximando la evolución de la fracción de vacíos como lineal cerca de la transición, $f_v(z) \approx f_{v,c} + k(z_c - z)$, y asumiendo que la función de escala $\mathcal{F}(x)$ para percolación en 3D se comporta asintóticamente como una sigmoide (perteneciente a la clase de universalidad de Ising/Percolación), obtenemos la forma analítica usada en el código:
\begin{equation}
\Theta(z) \approx \frac{1}{1 + \exp\left(\frac{z - z_c}{\Delta z}\right)}.
\end{equation}

El ancho de la transición $\Delta z$ no es arbitrario; está dictado por el exponente crítico $\nu \approx 0.88$ y la escala característica de la red $\xi_0$:
\begin{equation}
\Delta z \sim \left(\frac{\xi_0}{R_H(z_c)}\right)^{1/\nu}.
\end{equation}

Sustituyendo los valores optimizados en el análisis final, **$\Delta z \approx 0.014$** y **$z_c \approx 0.017$**, obtenemos una longitud de correlación implícita consistente con una escala de **$\sim 70$ Mpc**. Esto demuestra que la suavidad de la transición observada corresponde físicamente a la saturación del volumen de la **Hoja Local extendida**, respetando el límite superior cinemático impuesto por los datos de \textit{CosmicFlows-4}.

\section{Implementación Numérica y Reproducibilidad}
\label{app:codigo}

Este apéndice documenta la metodología numérica utilizada para confrontar la teoría GCAT con los datos de precisión de la era DESI/Euclid. El código completo incorpora priors bayesianos físicos derivados de la cinemática local para restringir la optimización de la transición topológica.

\subsection{B.1 Algoritmo de Optimización Global con Restricciones}

Dado que la superficie de verosimilitud ($\chi^2$) presenta mínimos locales espurios, utilizamos **Evolución Diferencial** (DE) modificada con funciones de penalización (barreras). El algoritmo busca el equilibrio entre los datos de expansión (que favorecen vacíos grandes) y los datos cinemáticos (que limitan el tamaño del vacío), convergiendo a la solución de saturación física en $z_c \sim 0.017$.

\subsubsection{Fórmulas Físicas Implementadas}

Las ecuaciones maestras implementadas en el módulo físico (\texttt{gcat\_physics.py}) imponen las siguientes condiciones:

\begin{enumerate}
    \item \textbf{Restricción de Información:} El factor de aliasing es fijo:
    \[
    \texttt{R\_FUND} = \frac{1}{6 \log_2 3} \approx 0.105155
    \]
    
    \item \textbf{Activación por Percolación (FSS):} La función de activación permite transiciones suavizadas ($\Delta z \sim 0.014$), reflejando que la percolación ocurre sobre el volumen extendido de la Hoja Local ($\sim 70$ Mpc) y no como un salto instantáneo:
    \[
    \Theta(z) = \frac{A_{\max}}{1 + \exp\left(\frac{z - z_c}{\Delta z}\right)}
    \]
    
    \item \textbf{Historia de Expansión Modificada:}
    \[
    H_{\text{GCAT}}(z) = H_{\Lambda\text{CDM}}(z) \cdot \left(1 - 2 \cdot \texttt{R\_FUND} \cdot \Theta(z)\right)^{-1/2}
    \]
\end{enumerate}

\subsection{B.2 Estructura del Código de Validación (V8)}

Presentamos el pseudocódigo de la función de coste final, que introduce explícitamente la restricción de \textit{CosmicFlows-4} como una barrera de penalización.

\begin{algorithmic}[1]
\REQUIRE Datos: $H(z)$ (incl. Pantheon+ bin local $z=0.015$), $S_8$, $H_0$
\FUNCTION{ChiSquaredTotal}{$\vec{\theta}$}
    \STATE \COMMENT{1. Calcular Dinámica de Percolación}
    \STATE $z_c, \Delta z \gets \vec{\theta}$
    \STATE $D_c \gets \text{ComovingDistance}(z_c)$
    
    \STATE \COMMENT{2. Aplicar Prior Físico (CosmicFlows-4)}
    \IF{$D_c > 70.0 \text{ Mpc}$}
        \STATE \RETURN $(D_c - 70.0) \times 100$ \COMMENT{Penalización de Muro (Mazurenko et al.)}
    \ENDIF
    
    \STATE \COMMENT{3. Evaluar Tensión de Hubble}
    \STATE $\Theta(z) \gets \text{SigmoidFSS}(z, z_c, \Delta z)$
    \STATE $H_{\text{pred}} \gets \text{FriedmannGCAT}(H_0^{\text{Planck}}, \Theta(0))$
    \STATE $\chi^2_{H_0} \gets ((H_{\text{pred}} - H_0^{\text{SH0ES}})/\sigma_{H_0})^2$
    
    \STATE \COMMENT{4. Evaluar Tensión S8}
    \STATE $S_8^{\text{pred}} \gets S_8^{\text{Planck}} \cdot (1 - \beta \cdot R_{\text{fund}} \cdot \Theta(0))$
    \STATE $\chi^2_{S_8} \gets ((S_8^{\text{pred}} - S_8^{\text{obs}})/\sigma_{S_8})^2$
    
    \STATE \COMMENT{5. Ajuste del "Escalón" Local y BAO}
    \STATE $\chi^2_{\text{Data}} \gets \sum ((H_{\text{GCAT}}(z_i) - H_{\text{obs}}(z_i))/\sigma_i)^2$
    
    \RETURN $\chi^2_{H_0} + \chi^2_{S_8} + \chi^2_{\text{Data}}$
\ENDFUNCTION
\end{algorithmic}

\subsection{B.3 Módulo de Ondas Gravitacionales}

El módulo de predicción se ha recalibrado con la escala de saturación identificada por la optimización.

\begin{algorithmic}[1]
\FUNCTION{GWResonance}{$\lambda_{\text{GW}}, z$}
    \STATE $\lambda_{\text{void}} \gets 70.0$ \COMMENT{Límite superior cinemático [Mpc]}
    \STATE \COMMENT{Resonancia estocástica centrada en $\sim 1.4 \times 10^{-16}$ Hz}
    \STATE $\mathcal{R} \gets \exp\left(-\frac{(\log \lambda_{\text{GW}} - \log \lambda_{\text{void}})^2}{2\sigma^2}\right)$
    \STATE $\text{Damping} \gets \exp\left(-\int_0^z R_{\text{fund}} \Theta(z') \mathcal{R} \frac{dz'}{H(z')}\right)$
    \RETURN $\text{Damping}$
\ENDFUNCTION
\end{algorithmic}

\subsection{B.4 Disponibilidad y Ciencia Abierta}

El repositorio incluye los scripts de optimización con restricciones (\texttt{scipy.optimize.differential\_evolution}), permitiendo a la comunidad verificar que la solución de 70 Mpc es un atractor robusto del sistema bajo las condiciones de \textit{CosmicFlows-4}.

Repositorio oficial: \url{https://github.com/[usuario]/GCAT-Cosmology-Validation}


% ============================================
% BIBLIOGRAFÍA ACTUALIZADA Y FINAL
% ============================================
\begin{thebibliography}{99}
\small

% ===== OBSERVACIONES CLAVE (ERA DE PRECISIÓN) =====
\bibitem{Planck2018}
Planck Collaboration.
\newblock Planck 2018 results. VI. Cosmological parameters.
\newblock \textit{Astronomy \& Astrophysics}, 641, A6, 2020.

\bibitem{Riess2022}
Riess, A. G., et al.
\newblock A Comprehensive Measurement of the Local Value of the Hubble Constant with 1 km/s/Mpc Uncertainty from the Hubble Space Telescope and the SH0ES Team.
\newblock \textit{The Astrophysical Journal Letters}, 934(1), L7, 2022.

\bibitem{DESI2024}
DESI Collaboration (Adame, A. G., et al.).
\newblock DESI 2024 VI: Cosmological Constraints from the Measurements of Baryon Acoustic Oscillations.
\newblock \textit{arXiv preprint}, arXiv:2404.03002, 2024.

\bibitem{Brout2022}
Brout, D., et al.
\newblock The Pantheon+ Analysis: Cosmological Constraints.
\newblock \textit{The Astrophysical Journal}, 938, 110, 2022.

\bibitem{KiDS2021}
Heymans, C., et al. (KiDS Collaboration).
\newblock KiDS-1000 Cosmology: Cosmic shear constraints and comparison between two point statistics.
\newblock \textit{Astronomy \& Astrophysics}, 646, A140, 2021.

\bibitem{DES2022}
Dark Energy Survey Collaboration.
\newblock Dark Energy Survey Year 3 results: Cosmological constraints from galaxy clustering and weak lensing.
\newblock \textit{Physical Review D}, 105(2), 023520, 2022.

% ===== ESTRUCTURA LOCAL Y FLUJOS PECULIARES (NUEVO) =====
\bibitem{Tully2023}
Tully, R. B., et al.
\newblock Cosmicflows-4.
\newblock \textit{The Astrophysical Journal}, 944(1), 94, 2023.

\bibitem{Mazurenko2024}
Mazurenko, N., et al.
\newblock Prior-free cosmological parameter estimation of Cosmicflows-4.
\newblock \textit{Monthly Notices of the Royal Astronomical Society}, 530(2), 2024. (Referencia clave para la restricción del Vacío Local).

% ===== FUNDAMENTOS TEÓRICOS (GNC & INFORMACIÓN) =====
\bibitem{Connes2008}
Chamseddine, A. H., \& Connes, A.
\newblock Why the Standard Model.
\newblock \textit{Journal of Geometry and Physics}, 58(1), 38–47, 2008.

\bibitem{Chamseddine2012}
Chamseddine, A. H., Connes, A., \& Marcolli, M.
\newblock Gravity and the standard model with neutrino mixing.
\newblock \textit{Advances in Theoretical and Mathematical Physics}, 11(6), 991–1089, 2007.

\bibitem{Hayes2001}
Hayes, B.
\newblock Third Base.
\newblock \textit{American Scientist}, 89(6), 490–494, 2001.

\bibitem{Palti2019}
Palti, E.
\newblock The Swampland: Introduction and Review.
\newblock \textit{Fortschritte der Physik}, 67, 1900037, 2019.

% ===== DINÁMICA Y FORMALISMO =====
\bibitem{Herglotz1930}
Herglotz, G.
\newblock Über die kanonischen Transformationen in der Variationsrechnung.
\newblock \textit{Anzeiger der Akademie der Wissenschaften in Wien}, 67, 103–107, 1930.

\bibitem{Lazo2017}
Lazo, M. J., et al.
\newblock Action principle for action-dependent Lagrangians: Toward nonconservative gravity.
\newblock \textit{Physical Review D}, 95(10), 101501, 2017.

% ===== PERCOLACIÓN Y MÉTODOS =====
\bibitem{Percolation3D}
Stauffer, D., \& Aharony, A.
\newblock \textit{Introduction to Percolation Theory}, 2nd ed.
\newblock Taylor \& Francis, London, 1994.

\bibitem{FSS}
Cardy, J. L.
\newblock \textit{Scaling and Renormalization in Statistical Physics}.
\newblock Cambridge University Press, Cambridge, 1996.

\bibitem{Storn1997}
Storn, R., \& Price, K.
\newblock Differential Evolution – A Simple and Efficient Heuristic for global Optimization over Continuous Spaces.
\newblock \textit{Journal of Global Optimization}, 11(4), 341–359, 1997.

\end{thebibliography}
\end{document}






